\subsection{Diagonal Matrices}

\begin{definitionbox}[Diagonal matrix]
A square matrix $D=(d_{ij})$ is diagonal when
\[
d_{ij}=0
\qquad\text{for every }i\neq j.
\]
We write
\[
D=\operatorname{diag}(d_1,d_2,\ldots,d_n).
\]
\end{definitionbox}

\begin{interpretationbox}[No mixing between coordinates]
A diagonal matrix acts independently on the coordinate positions. Since every
off-diagonal entry is zero, one coordinate does not contribute to another. This
is why products and powers of diagonal matrices are especially transparent.
\end{interpretationbox}

\begin{examplebox}[Powers of a diagonal matrix]
For
\[
D=\operatorname{diag}(3,-1,2),
\]
we have
\[
D^2=\operatorname{diag}(9,1,4),
\qquad
D^3=\operatorname{diag}(27,-1,8).
\]
Each diagonal entry is raised to the corresponding power.
\end{examplebox}

\begin{theorembox}[Diagonal matrices commute]
If
\[
D_1=\operatorname{diag}(a_1,\ldots,a_n),
\qquad
D_2=\operatorname{diag}(b_1,\ldots,b_n),
\]
then
\[
D_1D_2=D_2D_1
=\operatorname{diag}(a_1b_1,\ldots,a_nb_n).
\]
The commutativity follows from ordinary multiplication of the diagonal entries.
\end{theorembox}
